% --- Содержимое файла materials/5_1.tex ---
\subsection{Характеристическое свойство евклидовых пространств. Ба-
	наховы и гильбертовы пространства}
\begin{definition}[Евклидово пространство]
	\textbf{Евклидовым пространством} над полем $\mathbb{K}$ (где $\mathbb{K} = \mathbb{R}$ или $\mathbb{K} = \mathbb{C}$) называется линейное пространство $E$ над $\mathbb{K}$ с функцией $(\cdot, \cdot) : E^2 \to \mathbb{K}$, называемой \textbf{скалярным произведением}, которая обладает следующими свойствами:
	\begin{enumerate}
		\item $\forall x \in E: (x, x) \geqslant 0$, причем $(x, x) = 0 \Leftrightarrow x = 0$ (положительная определенность);
		\item $\forall x, y \in E: (x, y) = \overline{(y, x)}$ (эрмитова симметричность);
		\item $\forall x, y, z \in E, \forall \alpha, \beta \in \mathbb{K}: (\alpha x + \beta y, z) = \alpha(x, z) + \beta(y, z)$ (линейность по первому аргументу).
	\end{enumerate}
\end{definition}

\begin{theorem}[Характеристическое свойство Евклидовых пространств]
	Пусть $E$ --- линейное нормированное пространство. Норма в $E$ порождается скалярным произведением тогда и только тогда, когда для любых $x, y \in E$ выполняется \textbf{равенство параллелограмма}:
	\[\|x+y\|^2 + \|x-y\|^2 = 2\|x\|^2 + 2\|y\|^2\]
\end{theorem}


\begin{definition}[Типы пространств]~
	\begin{itemize}
		\item Полное линейное нормированное пространство $E$ называется \textbf{банаховым}.
		\item Полное евклидово пространство $H$ называется \textbf{гильбертовым}.
	\end{itemize}
\end{definition}

\begin{corollary}[Неравенство Коши---Буняковского]
	Пусть $E$ --- евклидово пространство. Тогда для любых $x, y \in E$:
	\[|(x, y)| \leqslant \|x\| \cdot \|y\|\]
\end{corollary}

\begin{corollary}
	Из неравенства Коши-Буняковского следует непрерывность скалярного произведения. Если $x_n \to x_0$ и $y_n \to y_0$, то:
	\[ |(x_n, y_n) - (x_0, y_0)| \leqslant |(x_n - x_0, y_n)| + |(x_0, y_n - y_0)| \leqslant \|x_n - x_0\|\|y_n\| + \|x_0\|\|y_n - y_0\| \to 0 \]
\end{corollary}